\documentclass{template}
\title{Installation Instructions for Python Code}

\author{Kyle Hansen}
\date{NE 795 Computational Project\\
14 December 2025}

\usepackage{cancel}
\usepackage{xcolor}


\usepackage{algpseudocode}
\usepackage{algorithm}

\usepackage{dirtree}

\newcommand*\Let[2]{\State{#1 $\gets$ #2}}
\algrenewcommand\algorithmicrequire{\textbf{Precondition:}}
\algrenewcommand\algorithmicensure{\textbf{Postcondition:}}


\newcommand*\multbox[1]{\fbox{\hspace{0ex}#1\hspace{0ex}}}

\newcommand{\I}{\mathcal{I}}
\newcommand{\E}{\mathcal{E}}

\definecolor{base03}{HTML}{002b36}
\definecolor{bg-light}{rgb}{0.95, 0.95, 0.95}

\newminted{python}{style=lightbulb, bgcolor=base03}
\newmintinline{python}{style=xcode, bgcolor=bg-light}

\begin{document}

\maketitle

\section{Instructions}

The code can be installed from the \href{https://github.com/hankyleh/Accelerated-Radiative-Diffusion#}{GitHub repository} using:

\begin{minted}[bgcolor=bg-light]{bash}
git clone https://github.com/hankyleh/Accelerated-Radiative-Diffusion.git
\end{minted}


The code requires Numpy, Scipy, and Matplotlib, and should be installed in an environment with these packages, or one should be created:

\begin{minted}[bgcolor=bg-light]{bash}
python -m venv venv
source venv/bin/activate
pip install numpy scipy matplotlib
\end{minted}


\section{Guide to operating code}

An input file should, at minimum, contain a function to compute group opacity, define $C_v$, initialize a variable of class Discretization, and define boundary and initial conditions. The method is called with

\end{document}